\section{Actual stopped mass and matching geometry}\label{sec:profile-proof}
The independent-report specialization below is the constrained-frontier mechanism already proved in~\cite{GTFStopped}; the comparison with correlated, controlled, model-indexed conditional rows is the extension. We give the argument in full because the history quantifiers and stopping-cost accounting are essential.
The reference measure in this section is a one-call acquired law, obtained from the reference preparation in Definition~\ref{def:recoverable}. It is not an assumption that any stopped policy has this law.

\begin{lemma}[Compact mass-constrained quantization]\label{lem:profile}
For every compact $K\subset\R^k$, probability $\nu$ supported on $K$, $m\ge1$ and real $t\ge1$, the minimum in \eqref{eq:profile} exists. For fixed centers put $d_C(x)=\min_j\norm{x-c_j}$. A minimizing density may be chosen in the form
\begin{equation}\label{eq:threshold-density}
 h(x)=t\one_{\{d_C(x)^2<\lambda\}}
       +t\xi\one_{\{d_C(x)^2=\lambda\}},\qquad 0\le\xi\le1,
\end{equation}
with $\int h\,d\nu=1$. The case $t=1$ permits $h=1$.
\end{lemma}
\begin{proof}
Centers may be projected onto the compact convex set $\conv K$, which cannot increase any distance to $K$. The set of densities $0\le h\le t$, $\int h=1$, is weak-star compact in $L^\infty(\nu)$. If centers converge, their squared nearest-center distance functions converge uniformly on $K$. The objective is therefore jointly continuous for this convergence and weak-star convergence of $h$. Compactness gives a minimum.

For fixed centers choose $\lambda$ such that
\[
 t\nu(d_C^2<\lambda)\le1\le t\nu(d_C^2\le\lambda)
\]
and choose $\xi$ to fill the possible boundary mass. Against any other admissible density $\widetilde h$, one has
\[
 \int(d_C^2-\lambda)(\widetilde h-h)\,d\nu\ge0,
\]
because on the sublevel set $h=t$ and on the strict superlevel set $h=0$. The equal total masses remove the $\lambda$ term. This proves optimality, including atoms and positive boundary mass. At a joint minimum the centers can also be replaced by the conditional means of their selected Voronoi cells, with any fixed Borel rule for ties; the squared-error projection identity proves that this does not increase the objective.
\end{proof}

\begin{lemma}[Acquisition mass bound]\label{lem:stopped-mass}
For any model-independent adaptive history policy and randomized stopping rule with $\E_\theta\tau\le N$, under the upper comparison in Definition~\ref{def:recoverable},
\begin{equation}\label{eq:actual-terminal-density}
 \law_\theta(X_\tau)\le b_*N\nu.
\end{equation}
More precisely, the unnormalized law of all acquired responses before and at stopping satisfies
\begin{equation}\label{eq:all-acquired}
 \E_\theta\sum_{t=1}^{\tau}\one_{\{X_t\in A\}}
       \le b_*\nu(A)\E_\theta\tau
       \quad(A\text{ Borel}).
\end{equation}
\end{lemma}
\begin{proof}
The event $\{\tau\ge t\}$ is determined before acquisition $t$, including all preceding private coins. Conditional on this information and the chosen action, the law of $X_t$ is at most $b_*\nu$. Summing
\[
 \E_\theta[\one_{\{\tau\ge t\}}\one_{\{X_t\in A\}}]
       \le b_*\nu(A)\Pp_\theta(\tau\ge t)
\]
over $t\ge1$ proves \eqref{eq:all-acquired} by Tonelli. The terminal indicator is no larger than the acquired-count sum. Its total mass is one, because finite mean implies almost-sure stopping, which proves \eqref{eq:actual-terminal-density}.
\end{proof}

\begin{lemma}[One-live-state realization]\label{lem:rejection}
Let $a_*N\ge1$. Given centers $C=\{c_1,\ldots,c_m\}$ and $0\le h\le a_*N$ with $\int h\,d\nu=1$, execute the fixed action $a_0$ and, after reading $x=f_{a_0}(y)$, accept with probability
\begin{equation}\label{eq:accept}
 A(x)=\frac{h(x)}{a_*N}.
\end{equation}
On acceptance, enter the nearest-center terminal label; otherwise return to the single live label. For every model this program satisfies
\begin{equation}\label{eq:rejection-bounds}
 \E_\theta\tau\le N,\qquad
 \E_\theta d_C(X_\tau)^2
       \le\frac{b_*}{a_*}\int d_C^2h\,d\nu.
\end{equation}
If all conditional response laws equal $\nu$, the choice $A=h/N$ with $0\le h\le N$ gives $\E\tau=N$ and $\law(X_\tau)=h\nu$ exactly.
\end{lemma}
\begin{proof}
At each surviving history the acceptance probability is at least
\[
 a_*\int\frac{h}{a_*N}\,d\nu=\frac1N.
\]
Consequently $\Pp_\theta(\tau>n)\le(1-1/N)^n$ and $\E_\theta\tau\le N$. The conditional expected squared distortion contributed by acceptance at that call is at most
\[
 \frac{b_*}{a_*N}\int d_C^2h\,d\nu.
\]
Summing it over surviving histories gives the second inequality of \eqref{eq:rejection-bounds}. This proof does not assume that repeated reports are independent or that the conditional law after rejection equals the original law.

For independent reference rows, the probability of acceptance at each surviving call is exactly $1/N$, and its accepted sublaw is $h\nu/N$. The sum of the geometric series is $N$, proving both exact assertions. Equivalently, the conditional response kernel being $\nu$ at every history suffices; no additional hidden-state independence is used.
\end{proof}

\begin{proof}[Proof of Theorem~\ref{thm:main}(ii)]
The protected audit identity gives, for every legal stopped policy and decoder,
\begin{equation}\label{eq:protected-pythagoras}
 \E_\theta\norm{Z-r_J}^2=B+\E_\theta\norm{X_\tau-r_J}^2.
\end{equation}
For a checkpoint there are at most $m$ decoder values. For an autonomous observer there are at most $M-1$, since at least one state is live. Conditional randomization of the decoder cannot improve squared loss; its mean belongs to $H$. Projecting centers onto $\conv K$ can only help. Lemma~\ref{lem:stopped-mass} bounds the actual terminal density by $b_*N$, so pointwise nearest-center comparison proves \eqref{eq:profile-check-lower}--\eqref{eq:profile-aut-lower} for every model and every randomized observer. It therefore proves the robust lower as well.

Apply Lemma~\ref{lem:rejection} to a minimizer in Lemma~\ref{lem:profile}, with $m=M-1$. This is a common program, determined by $\nu$, the comparison constants and its centers, not by $\theta$. It proves \eqref{eq:profile-aut-upper}. The same construction is available in the more permissive checkpoint architecture. If the conditional laws all equal $\nu$, the terminal law is exactly the minimizing $h\nu$, and the upper equals the lower. This proves \eqref{eq:exact-profile}. The stop is executed before the audit; no future audit information enters the accepting set.
\end{proof}

\begin{lemma}[Regular-support profile]\label{lem:regular}
Under \eqref{eq:regularity},
\begin{equation}\label{eq:profile-rate}
 c\,(tm)^{-2/s}\le\calQ_m^\nu(t)\le C\,(tm)^{-2/s}
 \quad(t\ge1,\ m\ge1),
\end{equation}
with constants depending only on the data specified in Theorem~\ref{thm:main}(iii).
\end{lemma}
\begin{proof}
For any center $c\in\conv K$, if $B(c,r)$ meets $K$, choose $x\in K\cap B(c,r)$. Then $K\cap B(c,r)\subset B(x,2r)$. Hence, when $2r\le r_0$, any union of $m$ such balls has $\nu$-mass at most $mC_0(2r)^s$. An admissible density is at most $t$. Choose a fixed sufficiently small $c_1>0$ and set $r=c_1(tm)^{-1/s}$, so that this acquired mass is at most $1/2$ and $2r\le r_0$. At least half the selected probability lies at distance at least $r$, proving the lower bound.

For the upper we record the covering and packing steps. A maximal $r$-separated subset of $K$ covers $K$ by $r$-balls, and its disjoint $r/3$-balls have mass at least $c_0(r/3)^s$. Thus its cardinality is at most $3^s/(c_0r^s)$. Enlarging constants to include the finitely small-cardinality regime gives a global cover by at most $m$ balls of radius $C_1m^{-1/s}$.

Conversely, choosing $\delta=c_2m^{-1/s}$ with $c_2$ sufficiently small, fewer than $m$ balls of radius $\delta$ cannot cover $K$, by the upper mass bound. A greedy construction therefore produces $m$ centers separated by at least $\delta$. For $t\ge t_0$, where $t_0$ depends only on the regularity constants, set
\[
 r=(c_0mt)^{-1/s}\le\delta/3.
\]
The $m$ disjoint $r$-balls have union $U$ of mass at least $mc_0r^s=1/t$. The density $h=\one_U/\nu(U)$ is admissible and its distance to these centers is at most $r$. This gives the required upper for $t\ge t_0$. For $1\le t<t_0$, use $h=1$ and the global $m$-center cover: $C_1^2m^{-2/s}\le C_1^2t_0^{2/s}(tm)^{-2/s}$. All changes for radii exceeding $r_0$ are absorbed by the fixed diameter constant, not by $t$ or $m$.
\end{proof}

\begin{proof}[Proof of Theorem~\ref{thm:main}(iii)]
For $a_*N\ge1$, insert Lemma~\ref{lem:regular} into the two profile bounds. For $1\le N<1/a_*$, stop after one report of $a_0$ and use the global $m$-center cover. This has mean acquisition one and excess at most $Cm^{-2/s}$, which is at most $Ca_*^{-2/s}(Nm)^{-2/s}$ in this range. The lower profile estimate holds for all $N\ge1$. The state deduction $m=M-1$ is exact throughout.
\end{proof}

\begin{proposition}[Atoms retain their physical mass]\label{prop:atoms}
For any compact acquired reference law and $m\ge1,t\ge1$,
\begin{equation}\label{eq:atom-threshold}
 \calQ_m^\nu(t)=0
 \quad\Longleftrightarrow\quad
 \text{some set of at most $m$ atoms has total $\nu$-mass at least $1/t$.}
\end{equation}
In particular, an atom of mass $w>0$ gives zero excess in the independent-reference experiment once $N\ge1/w$, with one live and one terminal state.
\end{proposition}
\begin{proof}
At a minimum with zero distortion, the probability measure $h\nu$ is supported on the at most $m$ centers. It can carry unit mass only if their $\nu$-mass is at least $1/t$. Conversely, restrict and normalize $\nu$ to such a finite set; the resulting density is at most $t$ and its distortion at those centers is zero. The causal assertion follows from the exact rejection implementation. All rejected observations are included in the mean $1/w$.
\end{proof}
